% Auto-generated complete chapter from 29C_artstation_breakdown_benchmark.md
\chapter{29C\_artstation\_breakdown\_benchmark}
\label{complete:root_039}

\begin{sourcemeta}
\textbf{Fuente:} \href{file:///E:/Laboral/29C_artstation_breakdown_benchmark.md}{29C\_artstation\_breakdown\_benchmark.md}\\
\textbf{Seccion:} root\\
\textbf{Estado:} reference\\
\textbf{Rol:} Benchmark o complemento\\
\textbf{Lineas originales:} 838\\
\textbf{Ultima modificacion:} 2026-06-11 18:27\\
\textbf{Deduplicacion conservadora:} 0 bloques repetidos omitidos; 0 caracteres repetidos no reimpresos.
\end{sourcemeta}

\section*{Resumen de orientacion}
The previous F3 output was \textbf{not complete enough} as a final \texttt{29C} module.

\section*{Contenido completo depurado}
\addcontentsline{toc}{section}{Contenido completo depurado}




\subsection{0. Module status}




\begin{mdcode}
Bloque de codigo: text
Module: 29C_artstation_breakdown_benchmark.md
Input task: F3_artstation_breakdown_benchmark
Date generated: 2026-06-11
Status: Complete v1
Source type: public portfolio-pattern benchmark + accessible snippet-derived examples
Primary use: ArtStation / portfolio breakdown structure for TwinSight X500 and Blender portrait
\end{mdcode}




\subsection{1. Completion audit}




The previous F3 output was \textbf{not complete enough} as a final \texttt{29C} module.




It had useful general guidance, but it lacked:




\begin{itemize}
\item a clear module structure;
\item source-quality classification;
\item separate treatment of TwinSight and the Blender portrait;
\item breakdown taxonomy;
\item ArtStation post architecture;
\item image requirements;
\item technical captions;
\item publication checklist;
\item mistakes to avoid;
\item direct update instructions for portfolio, GitHub and LinkedIn;
\item a retry prompt for future deeper ArtStation collection.
\end{itemize}




This version completes the module as an actionable benchmark and production guide.




\subsection{2. Scope}




This module defines how Alexander should present technical-art and real-time 3D breakdowns on:




\begin{itemize}
\item ArtStation;
\item personal portfolio;
\item GitHub README;
\item LinkedIn Featured;
\item interview follow-up materials.
\end{itemize}




It focuses on two assets:




\begin{enumerate}
\item \textbf{TwinSight X500}  
\end{enumerate}

   Unity WebGL technical visualization for drone assembly inspection.




\begin{enumerate}
\item \textbf{Blender hyperrealistic portrait}  
\end{enumerate}

   Supporting 3D art / character-art signal, useful only if presented with technical breakdown depth.




\subsection{3. Source limitations}




ArtStation often blocks automated access through security checks. Several pages and snippets were visible through search, but full profile/project inspection was not reliably available.




Therefore, this module should be treated as:




\begin{mdcode}
Bloque de codigo: text
Portfolio breakdown structure benchmark v1
\end{mdcode}




not as a complete database of verified ArtStation examples.




\subsection{4. Source quality table}




\begin{center}
\scriptsize
\begin{longtable}{>{\raggedright\arraybackslash}p{0.305\linewidth} >{\raggedright\arraybackslash}p{0.305\linewidth} >{\raggedright\arraybackslash}p{0.305\linewidth}}
\toprule
Source type & Reliability & Use \\
\midrule
\endfirsthead
\toprule
Source type & Reliability & Use \\
\midrule
\endhead
Fully accessible portfolio page & High & direct structure extraction \\
Search snippet from ArtStation & Medium & title/tag signal only \\
YouTube/Vimeo reel page & Medium & video structure \\
Studio case study & High & project narrative \\
General industry pattern & Medium & layout and content planning \\
Unopened ArtStation page & Low & do not cite as factual detail \\
\bottomrule
\end{longtable}
\end{center}




\subsection{5. Accessible benchmark signals}




\subsection{5.1 Search-snippet signals from ArtStation-like pages}




\begin{center}
\scriptsize
\begin{longtable}{>{\raggedright\arraybackslash}p{0.305\linewidth} >{\raggedright\arraybackslash}p{0.305\linewidth} >{\raggedright\arraybackslash}p{0.305\linewidth}}
\toprule
Signal & What it indicates & Use for Alexander \\
\midrule
\endfirsthead
\toprule
Signal & What it indicates & Use for Alexander \\
\midrule
\endhead
Technical Artist Demo Reel 2025 & reels are common for TA profiles & publish TwinSight video \\
Made with Unity & engine tag matters & tag Unity clearly \\
Unity foliage shader & shader breakdowns are portfolio assets & visual modes deserve breakdown \\
Unity URP water shader & URP is valid TA signal & state URP usage \\
Unity procedural fur shader & procedural/shader studies are valued & avoid claiming shader depth without proof \\
3D Rigger / Technical Artist Demo Reel & hybrid roles use reels & video must show systems \\
topology breakdown & expected for character/asset work & include wireframes \\
material breakdown & expected for realistic work & include texture/material maps \\
shader breakdown & expected for TA & include Shader Graph/code if present \\
game-ready breakdown & expected for real-time assets & include optimization metrics \\
\bottomrule
\end{longtable}
\end{center}




\subsection{5.2 Studio/case-study signals already collected}




\begin{center}
\scriptsize
\begin{longtable}{>{\raggedright\arraybackslash}p{0.455\linewidth} >{\raggedright\arraybackslash}p{0.455\linewidth}}
\toprule
Source & Relevant signal \\
\midrule
\endfirsthead
\toprule
Source & Relevant signal \\
\midrule
\endhead
HapzXR WebXR Product Configurator & browser-based 3D + product interaction \\
HapzXR Digital Twin Facility & spatial context + operational visualization \\
HapzXR BIM Overlay & model comparison + overlay logic \\
Virtual Verse Work Spatial & concise project page structure \\
Virtual Verse NBK Virtugate & WebGL + VR onboarding presentation \\
Industrial3D reels & industrial technical animation language \\
Kinetic Vision reel & technical visualization as product/process/system communication \\
\bottomrule
\end{longtable}
\end{center}




\subsection{6. Breakdown taxonomy}




\subsection{6.1 Technical art breakdown types}




\begin{center}
\scriptsize
\begin{longtable}{>{\raggedright\arraybackslash}p{0.305\linewidth} >{\raggedright\arraybackslash}p{0.305\linewidth} >{\raggedright\arraybackslash}p{0.305\linewidth}}
\toprule
Type & Typical content & Fit for TwinSight \\
\midrule
\endfirsthead
\toprule
Type & Typical content & Fit for TwinSight \\
\midrule
\endhead
Shader breakdown & graphs, nodes, render modes & High \\
Optimization breakdown & before/after geometry, texture budgets & Very high \\
Pipeline breakdown & CAD → Blender → Unity → WebGL & Very high \\
Tool breakdown & custom editor/runtime systems & High \\
Interaction breakdown & selection, UI, camera, states & Very high \\
Visual-mode breakdown & X-ray, ghosted, blueprint, thermal & Very high \\
Profiling breakdown & FPS, memory, load, draw calls & High \\
UX evaluation breakdown & SUS, NASA-TLX, Think-Aloud & Very high \\
Character art breakdown & topology, UVs, textures, grooming & Medium \\
Environment breakdown & scene assembly, lighting, props & Low-medium \\
\bottomrule
\end{longtable}
\end{center}




\subsection{6.2 Character / portrait breakdown types}




\begin{center}
\scriptsize
\begin{longtable}{>{\raggedright\arraybackslash}p{0.305\linewidth} >{\raggedright\arraybackslash}p{0.305\linewidth} >{\raggedright\arraybackslash}p{0.305\linewidth}}
\toprule
Type & Typical content & Fit for Blender portrait \\
\midrule
\endfirsthead
\toprule
Type & Typical content & Fit for Blender portrait \\
\midrule
\endhead
Sculpt breakdown & blockout to final sculpt & High \\
Topology breakdown & face loops, eyelids, mouth, deformation & High \\
UV breakdown & UDIMs/UVs, texture density & High \\
Skin material breakdown & albedo, roughness, SSS, displacement & High \\
Groom breakdown & hair guides, eyebrows, beard & High \\
Lighting breakdown & HDRI, key/fill/rim, camera & Medium-high \\
Render breakdown & passes, compositing & Medium \\
Real-time conversion & optimization for Unity/Unreal & High if done \\
\bottomrule
\end{longtable}
\end{center}




\subsection{7. ArtStation post architecture}




\subsection{7.1 Ideal structure for TwinSight ArtStation post}




\begin{mdcode}
Bloque de codigo: text
01. Cover image
02. 30–60 second embedded video
03. One-sentence project summary
04. Role / contribution
05. Final screenshots
06. Feature breakdown
07. CAD-to-realtime pipeline
08. Optimization metrics
09. Visual modes breakdown
10. Unity interaction systems
11. UI / technical panels
12. Evaluation results
13. Tools used
14. Limitations
15. Links
\end{mdcode}




\subsection{7.2 Ideal structure for Blender portrait ArtStation post}




\begin{mdcode}
Bloque de codigo: text
01. Cover render
02. Turntable video
03. Short project summary
04. Reference board
05. Sculpt stages
06. Topology
07. UVs
08. Texture maps
09. Skin shader / material setup
10. Grooming
11. Lighting
12. Final renders
13. Real-time/game-ready note if applicable
14. Tools used
15. Lessons learned
\end{mdcode}




\subsection{8. TwinSight breakdown specification}




\subsection{8.1 Cover}




The cover should not be a generic drone render.




Use one of:




\begin{center}
\scriptsize
\begin{longtable}{>{\raggedright\arraybackslash}p{0.455\linewidth} >{\raggedright\arraybackslash}p{0.455\linewidth}}
\toprule
Cover option & Strength \\
\midrule
\endfirsthead
\toprule
Cover option & Strength \\
\midrule
\endhead
Full drone with technical UI visible & best overall \\
Exploded view with labels & strongest assembly signal \\
Cross-section with UI panel & strongest inspection signal \\
Grid of visual modes & strongest technical-art signal \\
\bottomrule
\end{longtable}
\end{center}




Recommended cover title:




\begin{mdcode}
Bloque de codigo: text
TwinSight X500
Unity WebGL Technical Visualization
\end{mdcode}




Recommended subtitle:




\begin{mdcode}
Bloque de codigo: text
Drone assembly inspection · CAD-to-realtime · interactive 3D
\end{mdcode}




\subsection{8.2 Opening copy}




Use short copy:




\begin{mdcode}
Bloque de codigo: text
TwinSight X500 is a Unity WebGL technical visualization prototype for inspecting the assembly of a Holybro X500 V2 drone. It converts CAD-derived assets into an optimized browser-based 3D viewer with component selection, exploded view, clipping tools, visual modes and usability/workload evaluation.
\end{mdcode}




Do not open with:




\begin{mdcode}
Bloque de codigo: text
This was my thesis project about Design Science Research...
\end{mdcode}




Academic methodology should appear later.




\subsection{8.3 Role statement}




Use:




\begin{mdcode}
Bloque de codigo: text
Role:
Real-time 3D development, Unity WebGL integration, C# runtime systems, technical UI, visual modes, Blender optimization workflow, evaluation design and documentation.
\end{mdcode}




If AI assistance is mentioned:




\begin{mdcode}
Bloque de codigo: text
AI assistance was used for implementation support and iteration. Architecture, debugging, integration, validation, final decisions and project ownership remained my responsibility.
\end{mdcode}




\subsection{8.4 Key feature cards}




Use concise cards:




\begin{center}
\scriptsize
\begin{longtable}{>{\raggedright\arraybackslash}p{0.455\linewidth} >{\raggedright\arraybackslash}p{0.455\linewidth}}
\toprule
Feature & Caption \\
\midrule
\endfirsthead
\toprule
Feature & Caption \\
\midrule
\endhead
Component selection & Inspect individual drone parts and relationships \\
Exploded view & Reveal assembly hierarchy and spatial placement \\
Cross-section & Cut through structures to expose hidden geometry \\
Visual modes & Switch between realistic, X-ray, ghosted, blueprint and wireframe views \\
Technical UI & Display part information while inspecting the model \\
WebGL deployment & Run the viewer directly in a browser \\
Evaluation & Measure usability and workload using SUS and NASA-TLX \\
\bottomrule
\end{longtable}
\end{center}




\subsection{8.5 CAD-to-realtime section}




Recommended structure:




\begin{mdcode}
Bloque de codigo: text
CAD source
↓
conversion / tessellation
↓
Blender cleanup
↓
low-poly optimization
↓
UV / bake / material preparation
↓
Unity import
↓
WebGL deployment
\end{mdcode}




Include:




\begin{itemize}
\item before/after triangle counts;
\item wireframe comparison;
\item modular fastener system;
\item texture/material strategy;
\item constraints for mobile/web;
\item limitations.
\end{itemize}




\subsection{8.6 Optimization metrics section}




Use a metrics card layout.




\begin{center}
\scriptsize
\begin{longtable}{>{\raggedright\arraybackslash}p{0.455\linewidth} >{\raggedright\arraybackslash}p{0.455\linewidth}}
\toprule
Metric & Value \\
\midrule
\endfirsthead
\toprule
Metric & Value \\
\midrule
\endhead
Optimized geometry & 95,617 triangles \\
Source CAD routes & 6.5M+ triangles \\
Participants & 12 \\
SUS average & 91.88 \\
NASA-TLX Raw, 3D viewer & 8.69 \\
NASA-TLX Raw, 2D support & 19.89 \\
Task-condition records & 96 \\
\bottomrule
\end{longtable}
\end{center}




Add context:




\begin{mdcode}
Bloque de codigo: text
Evaluation was formative and academic, not a production deployment benchmark.
\end{mdcode}




\subsection{8.7 Visual modes breakdown}




Create a grid:




\begin{center}
\scriptsize
\begin{longtable}{>{\raggedright\arraybackslash}p{0.455\linewidth} >{\raggedright\arraybackslash}p{0.455\linewidth}}
\toprule
Mode & What to show \\
\midrule
\endfirsthead
\toprule
Mode & What to show \\
\midrule
\endhead
Realistic & final material view \\
X-ray & internal structure readability \\
Ghosted & partial transparency \\
Blueprint & technical drawing aesthetic \\
Wireframe & topology / engineering signal \\
Solid & neutral inspection view \\
Thermal-style & qualitative visual analysis, not real FEA \\
\bottomrule
\end{longtable}
\end{center}




Important wording:




\begin{mdcode}
Bloque de codigo: text
Thermal-style is a qualitative visual mode, not a physical thermal simulation.
\end{mdcode}




\subsection{8.8 Unity systems breakdown}




Include one compact section:




\begin{mdcode}
Bloque de codigo: text
Runtime systems:
- selection/highlighting
- camera control
- exploded-view state
- clipping plane control
- visual mode switching
- bottom-sheet UI
- part metadata panel
\end{mdcode}




Optional screenshot types:




\begin{itemize}
\item Unity hierarchy;
\item script inspector;
\item UI Toolkit panel;
\item Shader Graph;
\item C\# snippet;
\item profiler screenshot.
\end{itemize}




Do not include too much code on ArtStation. Put code detail in GitHub.




\subsection{9. Blender portrait breakdown specification}




\subsection{9.1 Purpose}




The Blender portrait should not be positioned as the main employability proof for technical visualization roles.




It should function as:




\begin{mdcode}
Bloque de codigo: text
supporting evidence of high-end 3D fundamentals, anatomy, material work, grooming and visual discipline.
\end{mdcode}




\subsection{9.2 Required sections}




\begin{center}
\scriptsize
\begin{longtable}{>{\raggedright\arraybackslash}p{0.455\linewidth} >{\raggedright\arraybackslash}p{0.455\linewidth}}
\toprule
Section & Required content \\
\midrule
\endfirsthead
\toprule
Section & Required content \\
\midrule
\endhead
Final render & 2–4 strongest images \\
Reference & board or brief reference explanation \\
Sculpt & blockout / mid / final \\
Topology & face loops and clean deformation \\
UVs & layout and density \\
Skin material & texture maps and shader setup \\
Eyes & iris/cornea material \\
Hair/groom & eyebrows, beard, eyelashes, hair guides \\
Lighting & setup diagram \\
Render & camera, focal length, compositing \\
Real-time adaptation & optional but valuable \\
\bottomrule
\end{longtable}
\end{center}




\subsection{9.3 What makes it relevant to technical art}




Add a section:




\begin{mdcode}
Bloque de codigo: text
Technical focus:
- topology decisions
- material layering
- groom setup
- render optimization
- texture organization
- possible real-time conversion path
\end{mdcode}




\subsection{9.4 What to avoid}




Do not frame it as:




\begin{mdcode}
Bloque de codigo: text
main portfolio proof for Unity WebGL roles.
\end{mdcode}




Do frame it as:




\begin{mdcode}
Bloque de codigo: text
supporting 3D craft and technical presentation.
\end{mdcode}




\subsection{10. Image requirements}




\subsection{10.1 TwinSight images}




\begin{center}
\scriptsize
\begin{longtable}{>{\raggedright\arraybackslash}p{0.455\linewidth} >{\raggedright\arraybackslash}p{0.455\linewidth}}
\toprule
Image & Priority \\
\midrule
\endfirsthead
\toprule
Image & Priority \\
\midrule
\endhead
Hero final viewer & Required \\
Exploded view & Required \\
Component selection & Required \\
Cross-section & Required \\
Visual modes grid & Required \\
CAD before/after & Required \\
Wireframe comparison & Required \\
UI panel close-up & Required \\
Evaluation metrics card & Required \\
Unity editor screenshot & Recommended \\
Blender cleanup screenshot & Recommended \\
WebGL browser screenshot & Required \\
GitHub / README preview & Optional \\
\bottomrule
\end{longtable}
\end{center}




\subsection{10.2 Blender portrait images}




\begin{center}
\scriptsize
\begin{longtable}{>{\raggedright\arraybackslash}p{0.455\linewidth} >{\raggedright\arraybackslash}p{0.455\linewidth}}
\toprule
Image & Priority \\
\midrule
\endfirsthead
\toprule
Image & Priority \\
\midrule
\endhead
Final portrait render & Required \\
Close-up face & Required \\
Side/profile view & Required \\
Wireframe overlay & Required \\
UV layout & Required \\
Texture maps & Required \\
Groom guides & Recommended \\
Lighting setup & Recommended \\
Clay render & Recommended \\
Turntable video & Recommended \\
\bottomrule
\end{longtable}
\end{center}




\subsection{11. Caption rules}




\subsection{11.1 Good caption style}




Use:




\begin{mdcode}
Bloque de codigo: text
CAD-derived mesh optimized for browser-based real-time inspection.
\end{mdcode}




Avoid:




\begin{mdcode}
Bloque de codigo: text
Here you can see the model that I worked on after importing it and making changes because it was too heavy and needed to work on the web.
\end{mdcode}




\subsection{11.2 TwinSight caption bank}




\begin{mdcode}
Bloque de codigo: text
Unity WebGL technical visualization prototype
Component-level drone assembly inspection
CAD-to-realtime optimization pipeline
Exploded view for spatial hierarchy
Cross-section tool for hidden structures
Technical UI with part information
Multiple visual modes for inspection
Browser-based deployment, no install required
SUS and NASA-TLX formative evaluation
\end{mdcode}




\subsection{11.3 Portrait caption bank}




\begin{mdcode}
Bloque de codigo: text
High-resolution sculpt
Clean facial topology
Skin material and texture layering
Eyebrow and beard grooming
Eye shader setup
Controlled studio lighting
Realistic render pass
Game-ready adaptation path
\end{mdcode}




\subsection{12. Tools section}




\subsection{12.1 TwinSight tools}




Use a clean tools list:




\begin{mdcode}
Bloque de codigo: text
Unity
C#
Unity WebGL
URP
UI Toolkit
Shader Graph
Blender
Marmoset Toolbag
RizomUV
GitHub
\end{mdcode}




Only include Marmoset/RizomUV if actually used in the final asset pipeline.




\subsection{12.2 Portrait tools}




\begin{mdcode}
Bloque de codigo: text
Blender
Cycles / Eevee Next
Substance 3D Painter
RizomUV / Blender UV tools
Marmoset Toolbag
Photoshop / Krita
\end{mdcode}




Only include actual tools.




\subsection{13. Tags}




\subsection{13.1 TwinSight ArtStation tags}




\begin{mdcode}
Bloque de codigo: text
Unity
Unity WebGL
Technical Art
Real-Time 3D
Technical Visualization
CAD Optimization
Drone
Interactive 3D
CSharp
URP
Blender
Digital Twin
Simulation
\end{mdcode}




Use \texttt{Digital Twin} cautiously.




\subsection{13.2 Blender portrait tags}




\begin{mdcode}
Bloque de codigo: text
Blender
Character Art
Portrait
Realistic Character
Topology
Grooming
Skin Shader
Lookdev
3D Art
Technical Art
\end{mdcode}




\subsection{14. Recommended layouts}




\subsection{14.1 TwinSight layout}




\begin{mdcode}
Bloque de codigo: text
Hero image
Video
Problem/Solution
Feature grid
Pipeline diagram
Optimization before/after
Interaction screenshots
Visual modes grid
Evaluation metrics
Tools
Links
\end{mdcode}




\subsection{14.2 Portrait layout}




\begin{mdcode}
Bloque de codigo: text
Hero render
Turntable
Reference
Sculpt process
Topology
UVs
Textures
Materials
Groom
Lighting
Final renders
Tools
\end{mdcode}




\subsection{15. ArtStation publication checklist}




\subsection{15.1 TwinSight checklist}




\begin{mdcode}
Bloque de codigo: text
[ ] Cover image clear at thumbnail size
[ ] Video under 90 seconds embedded near top
[ ] Project explained in first 2 sentences
[ ] Role clearly stated
[ ] Unity WebGL visible
[ ] Component selection shown
[ ] Exploded view shown
[ ] Cross-section shown
[ ] Visual modes shown
[ ] CAD optimization shown
[ ] Metrics shown
[ ] Limitations stated
[ ] GitHub link added
[ ] Portfolio link added
[ ] Demo link added
[ ] Tags added
\end{mdcode}




\subsection{15.2 Portrait checklist}




\begin{mdcode}
Bloque de codigo: text
[ ] Final render strong enough for cover
[ ] Turntable or image sequence included
[ ] References shown or described
[ ] Sculpt stages shown
[ ] Topology shown
[ ] UVs shown
[ ] Texture maps shown
[ ] Groom shown if relevant
[ ] Lighting shown
[ ] Tools listed
[ ] Technical contribution clear
\end{mdcode}




\subsection{16. Common mistakes}




\begin{center}
\scriptsize
\begin{longtable}{>{\raggedright\arraybackslash}p{0.455\linewidth} >{\raggedright\arraybackslash}p{0.455\linewidth}}
\toprule
Mistake & Why it hurts \\
\midrule
\endfirsthead
\toprule
Mistake & Why it hurts \\
\midrule
\endhead
Only showing final renders & hides technical skill \\
No wireframes & weakens real-time credibility \\
No metrics & weakens technical proof \\
Long academic intro & recruiters disengage \\
Too much text & ArtStation is visual-first \\
No role statement & ownership unclear \\
Unclear AI usage & trust risk \\
Broken demo links & credibility loss \\
Claiming digital twin too strongly & overclaim risk \\
Hiding WebGL constraints & weak technical honesty \\
\bottomrule
\end{longtable}
\end{center}




\subsection{17. Recommended TwinSight ArtStation copy}




\subsection{17.1 Short intro}




\begin{mdcode}
Bloque de codigo: text
TwinSight X500 is a Unity WebGL technical visualization prototype for drone assembly inspection. It transforms CAD-derived assets into an optimized browser-based 3D viewer with component selection, exploded view, cross-section tools, visual modes and technical UI panels.
\end{mdcode}




\subsection{17.2 Role}




\begin{mdcode}
Bloque de codigo: text
Role: Real-time 3D development, Unity WebGL integration, C# runtime systems, technical UI, visual modes, Blender optimization workflow, documentation and formative evaluation.
\end{mdcode}




\subsection{17.3 Limitations}




\begin{mdcode}
Bloque de codigo: text
This is an academic/portfolio prototype, not a deployed industrial digital twin. It does not include live IoT data, predictive maintenance, WebAR, or production monitoring. The thermal-style view is a qualitative visual mode, not a physical simulation.
\end{mdcode}




\subsection{17.4 Links section}




\begin{mdcode}
Bloque de codigo: text
Links:
- Live demo: [add link]
- GitHub: [add link]
- Case study: [add link]
- Demo video: [add link]
\end{mdcode}




\subsection{18. Recommended Blender portrait ArtStation copy}




\subsection{18.1 Short intro}




\begin{mdcode}
Bloque de codigo: text
A realistic Blender portrait study focused on facial anatomy, skin lookdev, grooming, lighting and technical presentation. The project is included as supporting evidence of 3D craft, material discipline and character-art fundamentals.
\end{mdcode}




\subsection{18.2 Role}




\begin{mdcode}
Bloque de codigo: text
Role: sculpting, modeling, material setup, grooming, lighting, rendering and breakdown documentation.
\end{mdcode}




\subsection{18.3 Technical focus}




\begin{mdcode}
Bloque de codigo: text
Technical focus:
- clean facial topology
- skin material layering
- eye material setup
- eyebrow/beard grooming
- lighting and render control
- possible real-time adaptation path
\end{mdcode}




\subsection{19. Integration with existing roadmap}




\subsection{19.1 Impact on role strategy}




This module does not change the main job-search strategy.




It reinforces:




\begin{mdcode}
Bloque de codigo: text
TwinSight must be the flagship.
The Blender portrait must be secondary.
ArtStation must show process, not just beauty shots.
\end{mdcode}




\subsection{19.2 Impact on portfolio order}




Recommended portfolio order:




\begin{mdcode}
Bloque de codigo: text
1. TwinSight X500
2. ARA Framework
3. Blender portrait
4. Smaller supporting projects
\end{mdcode}




ArtStation-specific order:




\begin{mdcode}
Bloque de codigo: text
1. TwinSight X500 technical breakdown
2. Blender portrait breakdown
3. Shader/visual-mode studies if separated
\end{mdcode}




\subsection{20. Updates required to other files}




\subsection{20.1 Update \texttt{20\_portfolio\_copy\_and\_site\_structure.md}}




Add:




\begin{mdcode}
Bloque de codigo: text
ArtStation publication flow:
- TwinSight technical breakdown
- portrait technical breakdown
- visual-mode study posts
\end{mdcode}




\subsection{20.2 Update \texttt{21\_twinsight\_demo\_video\_script\_and\_shotlist.md}}




Add:




\begin{mdcode}
Bloque de codigo: text
ArtStation-specific teaser export:
- 30–60 seconds
- no long intro
- burned-in captions
- feature-focused
\end{mdcode}




\subsection{20.3 Update \texttt{19\_github\_readme\_templates.md}}




Add:




\begin{mdcode}
Bloque de codigo: text
README should link to ArtStation breakdown for visual process.
ArtStation should link back to GitHub for code/process credibility.
\end{mdcode}




\subsection{21. Next recommended task}




Proceed with:




\begin{mdcode}
Bloque de codigo: text
30_live_job_postings_market_snapshot.md
\end{mdcode}




Recommended subtask:




\begin{mdcode}
Bloque de codigo: text
B1_unity_realtime3d_live_jobs
\end{mdcode}




\subsection{Tool recommendation}




\begin{center}
\scriptsize
\begin{longtable}{>{\raggedright\arraybackslash}p{0.455\linewidth} >{\raggedright\arraybackslash}p{0.455\linewidth}}
\toprule
Tool & Use \\
\midrule
\endfirsthead
\toprule
Tool & Use \\
\midrule
\endhead
Agent mode & Yes \\
Deep Research & No \\
\bottomrule
\end{longtable}
\end{center}




Reason:




The benchmark sequence for presentation artifacts now has:




\begin{mdcode}
Bloque de codigo: text
29 case-study benchmark
29B demo reels benchmark
29C ArtStation breakdown benchmark
\end{mdcode}




The next highest value is to collect live job postings and compare actual requirements against the prepared materials.




\subsection{22. Prompt for next task}




\begin{mdcode}
Bloque de codigo: text
Execute B1_unity_realtime3d_live_jobs.

Objective:
Collect live job postings for roles aligned with Alexander’s primary target:
Unity Developer, Real-Time 3D Developer, Unity Technical Artist, Unity WebGL Developer, Interactive 3D Developer, Technical Visualization Developer.

Candidate context:
Alexander is based in Colombia and targets remote international work. His strongest project is TwinSight X500, a Unity WebGL technical visualization prototype for drone assembly inspection with CAD-to-realtime optimization, C#, URP, UI Toolkit, Blender, exploded view, clipping, visual modes and usability/workload evaluation.

Output structured tables only.

Fields:
- Company
- Role title
- Job URL
- Country / region
- Remote scope
- Employment model
- Salary if visible
- Seniority
- Required skills
- Preferred skills
- Portfolio expectations
- GitHub expectations
- Work authorization notes
- Colombia remote eligibility
- Relevance to TwinSight
- Risk / mismatch
- Confidence
- Date checked

Rules:
- Use only live or recently listed postings.
- Do not assume remote means worldwide.
- Mark US-only, EU-only, LATAM-only, hybrid, onsite or unknown.
- Do not invent salary.
- If unavailable, write “not found”.
- Keep cells short.
- Include URLs.
- Tables only.
- No strategy.
\end{mdcode}



